
\section{Introduction}
Low-rank random feature models (RF-LR) are a simple yet remarkably effective way to leverage the benefits of wide neural networks while training only a small number of parameters. At each layer \( \ell \), the network linearly combines a frozen random nonlinear basis and trains only the output weights \( A^{(\ell)} \). This design brings three immediate advantages: (i) the trainable parameter budget scales linearly with width (\( O(rn) \) versus \( O(n^2) \) for dense layers), (ii) optimization avoids the Riemannian and curvature issues inherent to fully factorized training, and (iii) universal approximation is preserved because the random features maintain full support in function space. Equally important, in the kernel (NTK) regime the induced reproducing kernel Hilbert space (RKHS) is not weakened by depth: deep equals shallow for ReLU networks in the NTK sense, so RF-LR inherits the same RKHS as its fully deep counterparts while being far more parameter efficient~\cite{bietti2021deepequalsshallowrelu}.

The NTK viewpoint—initiated by~\cite{jacot2018neural} and extended to finite width and depth~\cite{arora2019exact,arora2019harnessing,yang2019scaling,terjek2022depth}—explains the lazy training regime as kernel regression with a data-dependent kernel fixed at initialization. Parallel to this, the random features literature shows how to approximate kernels efficiently with finite dictionaries~\cite{rahimi2007random,rahimi2008weighted}, together with generalization guarantees and sampling strategies~\cite{bach2017breaking,rudi2017generalization,sriperumbudur2015optimal}. RF-LR lies at the intersection: it retains an NTK description while realizing it with a small trainable top layer over a large, frozen dictionary. Empirically and theoretically, this often yields better conditioning and concentration (e.g., doubly-random Gram regimes with \( O(1/r^2) \) fluctuations), making the deterministic kernel limit accurate already at moderate widths.

Our focus is twofold. First, we formulate and analyze the NTK of RF-LR at initialization, deriving exact recursions and finite-width corrections (building on classical NNGP/NTK recursions). Second, we develop a spectral theory in the two-layer setting: a signal–noise decomposition for the empirical kernel and a deformed Marchenko–Pastur (MP) description of the spectrum with a rank-one spike. Together, these explain the staircase loss curves observed during training: the NTK spectrum is typically dominated by one outlier and a concentrated bulk; early training tracks the outlier direction, then progressively activates higher-frequency components, matching the spectral-bias narrative~\cite{rahaman2019spectral,xu2019training,basri2020frequency,tancik2020fourier}.

A central objective of this work is to reconcile the NTK framework with weight parsimony: we aim for kernel-level understanding while training as few parameters as possible. In parallel, prior work on Multi-Component and Multi-Layer Neural Networks (MMNN) reports super convergence phenomena and striking stepwise loss behavior~\cite{zhang2025structuredbalancedmulticomponentmultilayer}. We focus on explaining the first training phase in these observations—the early/lazy regime—through an RF-LR NTK lens: how the kernel spectrum and its leading outlier drive fast initial error decay before higher-frequency components are activated.

Related work is extensive and we synthesize it here as a cohesive landscape. The NTK framework of~\cite{jacot2018neural} established that infinitely wide networks trained by gradient descent behave like kernel methods, with numerous extensions to realistic depths and architectures~\cite{arora2019exact,arora2019harnessing,yang2019scaling,terjek2022depth}. Random features~\cite{rahimi2007random,rahimi2008weighted} provide a practical route to kernel approximation with strong generalization guarantees in high dimension~\cite{bach2017breaking,rudi2017generalization} and robust statistical control~\cite{sriperumbudur2015optimal}. In parallel, a rich line of work investigates spectral bias—why low frequencies are learned first and high frequencies later—both empirically and theoretically~\cite{rahaman2019spectral,xu2019training,basri2020frequency,tancik2020fourier}. Low-rank and factorized architectures~\cite{sainath2013low,novikov2015tensorizing,arora2018optimization} seek parameter efficiency and implicit regularization; RF-LR can be seen as an extreme in which feature weights are frozen and only linear heads are trained, enabling sharp concentration (e.g., \( O(1/r^2) \) in doubly-random Gram regimes) and clean NTK analysis. Beyond these, the optimization landscape and stability literature (e.g., edge-of-stability phenomena~\cite{cohen2021gradient}, the role of adaptive optimizers~\cite{kingma2014adam,wilson2017marginal}) informs practical training choices; and neural scaling laws~\cite{kaplan2020scaling,hoffmann2022training,wei2022emergent} motivate architectures (like RF-LR) that deliver competitive accuracy at controlled compute.

Technically, we adopt and leverage the “deep equals shallow” result for ReLU kernels in the NTK regime~\cite{bietti2021deepequalsshallowrelu}. Concretely, we use the theorem (without restating it) to justify that the RKHS induced by deep RF-LR coincides with that of a shallow ReLU kernel. Alongside this equivalence, the same work characterizes the eigenvalue decay of rotation-invariant kernels on the sphere via endpoint expansions of the kernel \( \kappa(\rho) \) near \( \rho=\pm 1 \): under \( C^\infty(-1,1) \) and non-integer exponent \( \nu>0 \), the spherical-harmonic eigenvalues satisfy \( \mu_k \sim C(d,\nu)\,k^{-d-2\nu+1} \) up to parity-dependent coefficients determined by the leading terms \( c_1,c_{-1} \) in the expansions; if the kernel is \( C^\infty \) up to the endpoints with no such \( \nu \), the decay is faster than any polynomial. We adopt these spectral rates as our baseline regularity model to connect RKHS smoothness, kernel spectra, and training behavior.

From a scaling perspective, RF-LR brings a favorable tradeoff. Achieving and maintaining NTK behavior in fully-trainable deep networks can impose stringent width–data requirements (often quoted in the literature as polynomial in the sample size \( N \), with worst-case exponents tied to depth and curvature). In contrast, because RF-LR trains only \( O(rn) \) parameters and concentrates quickly, the NTK regime can be reached at scales effectively proportional to \( N\cdot r \) rather than \( N^4 \) in pessimistic dense cases (we formalize this improvement in our scaling proposition). In short, RF-LR converts “depth-induced” width requirements into a linear rank budget while suffering no RKHS disadvantage.

Our contributions are therefore: (i) an NTK recursion specialized to RF-LR together with finite-width corrections and concentration rates, (ii) a two-layer spectral theory yielding a signal–noise decomposition of the empirical kernel and a deformed MP law with a spike for the kernel matrix spectrum, and (iii) a principled argument—grounded in~\cite{bietti2021deepequalsshallowrelu}—that RF-LR suffers no RKHS disadvantage relative to deep fully-trainable networks. These theoretical results are complemented by numerical experiments (omitted here) that verify NTK spectral concentration, staircase loss dynamics, and the scaling benefits of RF-LR at moderate widths.

\vspace{1em}
\subsection{Related work}

 
The Neural Tangent Kernel (NTK) framework~\cite{jacot2018neural} characterizes infinitely wide neural networks trained by gradient descent as kernel methods with a fixed kernel. Extensions to finite width~\cite{arora2019exact}, convolutional architectures~\cite{arora2019harnessing}, and depth~\cite{yang2019scaling,terjek2022depth} have been developed. Our work contributes explicit NTK recursions for low-rank random feature models with \( O(1/r) \) fluctuation bounds, extending spectral analysis to architectures with bottlenecks.

 
Random feature approximations~\cite{rahimi2007random} approximate kernel methods by sampling features from a distribution. Theoretical guarantees for generalization~\cite{bach2017breaking,rudi2017generalization} and approximation~\cite{sriperumbudur2015optimal} have been established. Our frozen random features maintain full support throughout training, providing universal approximation at all finite times—a crucial property for mean-field convergence that differs from standard kernel ridge regression.

% (removed mean-field theory paragraph)

 
Low-rank factorizations~\cite{sainath2013low,novikov2015tensorizing} reduce parameter counts in neural networks. Training dynamics of factorized weights~\cite{arora2018optimization} differ from unfactorized gradient descent due to implicit regularization. Our work distinguishes RF-LR (training output weights \( A^{(\ell)} \) only) from full factorization approaches (training both factors), avoiding Riemannian geometry complications while achieving linear parameter scaling \( O(rn) \).

 
Neural networks exhibit spectral bias toward low frequencies~\cite{rahaman2019spectral,xu2019training}, learning smooth functions before complex ones. This bias is characterized via NTK analysis~\cite{basri2020frequency} and Fourier analysis~\cite{tancik2020fourier}. We extend this to stepwise (staircase) loss curves via band-by-band Fourier mode activation in mean-field dynamics, connecting spectral bias to multi-index learning.

 
Multi-index models~\cite{benarous2021online,damian2022neural,abbe2023sgd} learn functions of the form \( f(x) = \sum_{j=1}^r g_j(\langle v_j, x \rangle) \) by discovering latent directions \( v_j \) sequentially. Training exhibits staircase dynamics with plateaus corresponding to inactive directions. Extensive-rank regimes \( r \asymp d^\beta \) yield scaling laws~\cite{damian2022neural}. We show RF-LR realizes multi-index learning in Fourier/spherical harmonic space via cosine similarity geometry, while going beyond it in 1D via direct spatial localization.

 
Dictionary learning~\cite{mairal2009online,elad2010sparse} seeks sparse representations in overcomplete bases. Neural networks perform implicit dictionary learning~\cite{papyan2017convolutional,sulam2020multi}. Our random feature dictionary \( \{\ReLU(\langle w_j, x \rangle)\}_{j=1}^n \) is overcomplete and fixed; training selects sparsely (rank \( r \ll n \)) which features to use, yielding wavelet-like localized functions.

 
Gradient descent dynamics near sharp minima exhibit instability~\cite{cohen2021gradient}. Edge-of-stability (EOS) phenomena~\cite{cohen2021gradient} explain loss oscillations when eigenvalues exceed \( 2/\eta \). Adaptive optimizers (Adam~\cite{kingma2014adam}) behave differently from SGD, with distinct implicit biases~\cite{wilson2017marginal}. We document Adam's instability in sharp regimes and advocate optimizer switching (Adam \( \to \) SGD) for RF-LR training.

 
Neural scaling laws~\cite{kaplan2020scaling,hoffmann2022training} describe test loss as power laws in compute, data, and model size. Emergent capabilities~\cite{wei2022emergent} appear at critical scales. Our stepwise loss curves exhibit similar emergent behavior, with each drop corresponding to activating new frequency bands, explained by extensive-rank multi-index theory.

 
Universal approximation theorems~\cite{cybenko1989approximation,hornik1989multilayer} guarantee that neural networks can approximate continuous functions. Random feature models also satisfy universal approximation~\cite{rahimi2008weighted} under sufficient width. Three-layer networks with frozen features maintain approximation power at \emph{all training times}~\cite{nguyen2021proof}, not just at convergence—a key property we exploit.

 
Loss landscape geometry~\cite{li2018visualizing,fort2019large} affects trainability. Mean-field analysis reveals landscape convexification during training~\cite{chizat2020implicit}. We quantify this via Hessian spectrum evolution, showing sharpness transitions at frequency band activations and explaining when SGD switching becomes beneficial.

\subsection*{Our Contributions}
\begin{itemize}
    \item \textbf{NTK theory with no RKHS disadvantage and linear width scaling.} We prove that RF-LR induces the same RKHS as standard fully-connected networks for large rank \( r \), ensuring no expressivity loss. Moreover, the NTK regime is achieved at \emph{linear} parameter budget \( O(rn) \) rather than quadratic \( O(n^2) \), dramatically reducing the width required to enter the kernel regime while maintaining full theoretical guarantees. This establishes low-rank random features as strictly more efficient than dense architectures for lazy training.
    \item \textbf{NTK recursion} We derive the NTK recursion expansion for RF-LR without bias/output scaling (\( \sigma_A, \sigma_c, \beta = 1 \)), yielding clean expressions for two-layer and general \( L \)-layer networks.
    \item \textbf{Low-rank fluctuation bounds with improved concentration.} We formalize how rank-\( r \) bottlenecks yield \( O(1/r) \) kernel fluctuation decay (or \( O(1/r^2) \) in doubly-random Gram matrix regimes). Crucially, the \( O(1/r^2) \) decay is \emph{exponentially faster} than the standard \( O(1/N) \) decay from classical NTK theory, establishing that low-rank random features improve both parameter efficiency and kernel approximation reliability at finite width.
    \item \textbf{Spectral theory (two-layer case).} We state and analyze Fisher–Kibble decoupling, a signal–noise decomposition of the empirical NTK, and a deformed Marchenko–Pastur law (with spike) for the kernel matrix spectrum under standard random matrix scaling. Proofs are deferred to the appendix.
\end{itemize}
